\documentclass[11pt,a4paper]{article}

\input{../../sablony/pripravy/priprava.sty}

\nastavPripravu{4. ročník}{Kombinatorika}{36}{Kombinace I -- základy a kombinační číslo}

\begin{document}

\blok{Motivace a návaznost}

Z devíti studentů vybíráme tříčlennou komisi bez funkcí. Proč nyní různé pořadí týchž tří osob neznamená jinou komisi?

\poznamka{Předpoklady: variace, permutace, faktoriál. Cíl: vysvětlit význam neuspořádaného výběru a proč dělíme počtem uspořádání vybraných prvků.}

\blok{Výklad a odvození}

\begin{definiceBox}
\textbf{Kombinace bez opakování} délky $k$ z $n$ různých prvků je neuspořádaný výběr $k$ různých prvků. Jejich počet udává \textbf{kombinační číslo}
\[
\binom nk=\frac{n!}{k!(n-k)!},\qquad 0\le k\le n.
\]
\end{definiceBox}
Každou $k$-prvkovou skupinu počítají variace právě $k!$-krát. Proto $\binom nk=V(n,k)/k!$. Platí $\binom n0=\binom nn=1$; prázdná skupina je jedna možnost. Vždy zkontroluj, zda záleží na pořadí a zda jsou objekty různé.

\blok{Společně řešené příklady}

\begin{prikladReseniBox}[title={Příklad 1}]
Kolik různých tříčlenných komisí vznikne výběrem z devíti studentů?
\tcblower
$\binom93=\frac{9\cdot8\cdot7}{3\cdot2\cdot1}=84$.
\end{prikladReseniBox}



\begin{prikladReseniBox}[title={Příklad 2}]
Kolik úseček určuje šest různých bodů?
\tcblower
Každá dvojice různých bodů určuje úsečku, tedy $\binom62=15$. Poloha ostatních bodů výsledek nemění.
\end{prikladReseniBox}



\begin{prikladReseniBox}[title={Příklad 3}]
Kolika způsoby lze z osmi knih vybrat pět, pokud nezáleží na pořadí?
\tcblower
$\binom85=\frac{8!}{5!3!}=56$.
\end{prikladReseniBox}



\blok{Samostatná práce}

\begin{zadaniReseniBox}[title={Samostatná práce 1}]
Urči $\binom73$.
\tcblower
$35$.
\end{zadaniReseniBox}



\begin{zadaniReseniBox}[title={Samostatná práce 2}]
Ze šesti žáků vybereme dvojčlenný tým bez funkcí. Kolik možností?
\tcblower
$\binom62=15$.
\end{zadaniReseniBox}



\begin{zadaniReseniBox}[title={Samostatná práce 3}]
Porovnej počet dvojic členů komise a dvojic předseda--místopředseda ze šesti osob.
\tcblower
Komise: $\binom62=15$; funkce: $6\cdot5=30$.
\end{zadaniReseniBox}



\blok{Rozšíření}

\begin{zadaniReseniBox}[title={Rozšíření 1}]
Urči $\binom{12}{10}$ bez výpočtu velkých faktoriálů.
\tcblower
$\binom{12}{10}=\frac{12\cdot11}{2}=66$.
\end{zadaniReseniBox}



\blok{Domácí úkol}

\poznamka{Doporučená práce doma: přibližně 15--30 minut. Úlohy jsou seřazeny od základních k náročnějším; učitel může poslední podúlohu označit jako volitelnou.}

\begin{zadaniReseniBox}[title={Domácí úkol}]

\begin{enumerate}[label=\alph*),leftmargin=*,itemsep=2pt]

\item Vypočti $\binom82$ a $\binom84$.
\item Ze sedmi osob vyber tříčlennou komisi bez funkcí.
\item Kolik různých čtyřprvkových podmnožin má osmiprvková množina?
\item Vypočti $\binom{10}{8}$ zkrácením faktoriálů.
\item Vysvětli, proč při výběru předsedy a zástupce ze šesti lidí nepoužijeme přímo $\binom62$.
\item Z deseti osob vyber čtyřčlenný tým, v němž musí být Anna. Kolik možností?

\end{enumerate}

\tcblower

\begin{enumerate}[label=\alph*),leftmargin=*,itemsep=2pt]

\item $28$, resp. $70$.
\item $\binom73=35$.
\item $\binom84=70$.
\item $45$.
\item Role jsou rozlišené, a proto jde o $V(6,2)=30$, nikoli kombinaci.
\item Anna je pevná; tři další osoby vybíráme z devíti: $\binom93=84$.

\end{enumerate}

\end{zadaniReseniBox}

\poznamka{Orientační práce v hodině: 0--12 min porovnání s variacemi; 12--25 min odvození; 25--40 min úlohy; 40--45 min shrnutí.}

\end{document}
